\chapter[Frucht und Segen der Kirchbaukassen]{Ein Bericht über Frucht und Segen unserer Distrikts-Kirchbaukassen}
\quellseite{6}

Lieber Pastor H—! Ihren lieben Brief habe ich letzte Woche empfangen. Da unsere Vorsteher noch fleißig am Collectiren waren für die \$200.00, so antworte ich erst jetzt. Schicke Ihnen zugleich die letzten \$200.00 für die Kirchbaukasse im Namen der Gemeinde mit vielem und herzlichem Dank zurück. Gott möge es den Beamten der Kirchbaukasse und den vielen Gebern vergelten, was sie bei der ersten und auch bei der zweiten Anleihe Gutes an uns gethan haben. Zugleich freue ich mich auch herzlich, daß wir jedes Jahr im Stande gewesen sind, den versprochenen Theil der geborgten Summe zurückzubezahlen. Ich muß sagen, gerade bei dem von der Kirchbaukasse geborgten Gelde hat meine Gemeinde gelernt, Schulden abzutragen. Denn das stand ihr von vornherein fest: dieses Geld muß zur bestimmten Zeit zurückbezahlt werden, und wenn wir zweimal im Jahr unter unsern Gliedern collectiren müssen. So ist die Kirchbaukasse wirklich ein Segen für eine junge Gemeinde, bei der sie lernt zurückzugeben, was geborgt ist. Sodann hilft die Kirchbaukasse überhaupt im Allgemeinen dazu, daß junge Gemeinden das Geben lernen. Ich will meine Gemeinde nicht vor andern rühmen. Ich denke manchmal: Viele deiner Leute könnten fleißiger sein im Geben, aber wenn ich dann wieder das berechne, wie viel jährlich auf das stimmberechtigte Glied kommt, um unsere laufenden Ausgaben zu bezahlen und Schulden abzutragen, so finde ich viel Ursache, Gott zu loben und ihm zu danken. Nicht um meine Gemeinde zu rühmen, sondern zum Preise Gottes, und um den Beamten der Kirchbaukasse Muth zu machen, im Vertrauen auf Gott für diese Sache weiter zu reden und zu schreiben, möge hier zum Schluß noch ein kurzes Rechenexempel folgen. Meine Gemeinde, die jetzt aus 70 Gliedern besteht — letztes Jahr waren es 65 —, und zwar zum Theil aus recht armen Leuten, hat im letzten Jahre \$1750.00 für ihren eigenen Haushalt aufgebracht und dann noch etwa \$400.00 für Schuldentilgung. Das macht also für jedes stimmberechtigte Glied jährlich etwas über \$30.00. Ich glaube, in wenig Gemeinden ist die Durchschnittssumme so hoch wie bei uns. Daneben hat meine Gemeinde für wohltätige Zwecke nach außen hin noch \$181.22 aufgebracht. Ich habe dieses deshalb mitgetheilt, um den Beamten der Kirchbaukasse zu zeigen, daß es sich wirklich lohnt, wenn diese Kasse jungen und hilfsbedürftigen Gemeinden in der Zeit ihrer Noth mit einer zinsfreien Anleihe unter die Arme greift. Gott segne daher auch weiter das Missionswerk der Kirchbaukasse und helfe, daß meine nun erstarkte Gemeinde sich auch in Zukunft recht dankbar erweise für diese Wohlthat. Mit herzlichem Gruße und Dank

\begin{unterschrift}
E. H., den 26. October 1903.\\
Ihr E. S.
\end{unterschrift}

Und die Verwalter obgenannter Kasse bitten um ein Räumlein in den Spalten des „Lutheraner“ für obige Zeilen, weil sie alle Christen, die für die Sache des Reiches ihres lieben Heilandes und besonders auch für diese unter uns noch nicht alte Weise der Beihülfe zur Ausbreitung desselben eifrig thätig sind, erfreuen und anderen das Herz dafür warm machen werden. Ein solches Beispiel überzeugt mehr als viele Worte. Und doch, dürfen wir noch einige Worte anfügen? Ein Gemeindlein ist gesammelt. Nicht selten haben sich einige wenige Christen zusammengepaart. Unter sich haben sie keine Millionäre. Es fehlt ein Kirch- oder Schulgebäude. Zwar tragen sie selbst zusammen, dürfen auch wohl Gaben von Schwestergemeinden in ihre Kasse thun, aber es gibt Schulden. Zinsen stecken einen Löffel mit in die Suppe. Zumal eine junge Gemeinde trägt schwer an einer Zinsenlast. Ihre Ausbreitung wird gehemmt. Furcht vor den Zinsen verhindert nicht selten die Errichtung nötiger Bauten. Und selbst solche Anleihen müssen erst gesucht werden und sind schwer zu erlangen. Um solcher Noth abzuhelfen, sind die Gemeinden des Distrikts zusammengetreten, haben zusammengelegt und bringen immer noch dazu. Das ist die sogenannte Kirchbaukasse. Daraus wird zinsfrei geliehen; sobald ein Sümmlein heimkehrt, wird die Hand eines andern Gemeindleins gefüllt. Manches Kirchlein, nicht wenige Schulgebäude, hie und da ein Pfarrhaus hat so erstehen können. Und dann kommt der Diener, das geliehene Geld, wieder zurück, und wieder ein Kirchlein oder ein Schulgebäude oder auch ein Pfarrhaus wird mit Hülfe desselben Geldes anderswo erbaut. Und die Gemeinde lernt auf eigenen Füßen stehen, ja, hilft aus Dankbarkeit andern Häuslein, wie obiger Brief darthut. „Und der Herr, unser Gott, sei uns freundlich und fördere das Werk unserer Hände bei uns, ja, das Werk unserer Hände wolle er fördern.“ Im Auftrage L. Hölter.
